% ============================================================================
%  PLAN DE CALIDAD DE SOFTWARE — ChibchaWeb
%  Documento compilable con pdflatex o lualatex
% ============================================================================
\documentclass[12pt,a4paper]{article}

% — Codificación y tipografía ------------------------------------------------
\usepackage[utf8]{inputenc}
\usepackage[T1]{fontenc}
\usepackage[spanish]{babel}
\usepackage{lmodern}

% — Diseño de página ---------------------------------------------------------
\usepackage[margin=2.5cm]{geometry}
\usepackage{parskip}

% — Tablas y figuras ----------------------------------------------------------
\usepackage{booktabs}
\usepackage{tabularx}
\usepackage{longtable}
\usepackage{multirow}
\usepackage{array}

% — Colores y resaltado -------------------------------------------------------
\usepackage[table,dvipsnames]{xcolor}
\definecolor{secbg}{HTML}{EBF5FB}

% — Hipervínculos y metadatos -------------------------------------------------
\usepackage{hyperref}
\hypersetup{
  colorlinks=true,
  linkcolor=MidnightBlue,
  urlcolor=MidnightBlue,
  pdfauthor={Equipo ChibchaWeb},
  pdftitle={Plan de Calidad de Software — ChibchaWeb},
}

% — Encabezados y pies de página ----------------------------------------------
\usepackage{fancyhdr}
\pagestyle{fancy}
\fancyhf{}
\fancyhead[L]{\small\textsc{ChibchaWeb}}
\fancyhead[R]{\small Plan de Calidad de Software}
\fancyfoot[C]{\thepage}
\renewcommand{\headrulewidth}{0.4pt}

% — Enumeración personalizada -------------------------------------------------
\usepackage{enumitem}

% ============================================================================
\begin{document}

% — Portada -------------------------------------------------------------------
\begin{titlepage}
\centering
\vspace*{2cm}
{\Huge\bfseries Plan de Calidad de Software\\[0.4em]
Proyecto ChibchaWeb\par}
\vspace{1.5cm}
{\Large\itshape Definición de procesos, métricas, responsabilidades\\y herramientas de aseguramiento de calidad\par}
\vspace{2cm}
{\large
\begin{tabular}{rl}
\textbf{Proyecto:}   & ChibchaWeb \\
\textbf{Fecha:}      & 26 de mayo de 2026 \\
\end{tabular}\par}
\vspace{1.5cm}

{\large\textbf{Integrantes del equipo}\par}
\vspace{0.5cm}
\begin{tabular}{lll}
\toprule
\textbf{Nombre y apellidos} & \textbf{Código} & \textbf{Rol en el proyecto} \\
\midrule
Juan Carlos Duarte Sandoval & 20212020149 & Desarrollador \\
Sebastián Alejandro Bernate Guzmán & 20211020121 & Desarrollador \\
Laura Carina Álvarez Campos & 20212020006 & Product Owner \\
Johan Andrés Tamara Flautero & 20211020039 & Desarrollador \\
Laura Daniela Cubillos Escobar & 20211020045 & Product Manager \\
David Santiago Torres Mendieta & 20211020144 & Analista QA \\
\bottomrule
\end{tabular}

\vfill
{\small Universidad Distrital Francisco José de Caldas}
\end{titlepage}

% — Índice --------------------------------------------------------------------
\tableofcontents
\newpage

% ============================================================================
\section{Introducción}
% ============================================================================

ChibchaWeb es una plataforma profesional de hosting web con presencia destacada en
Colombia. Ofrece servidores SSD NVMe de alto rendimiento, disponibilidad garantizada del
99,9\,\%, certificados SSL gratuitos, respaldos automáticos y soporte técnico especializado
disponible las 24 horas del día.

Este documento, el Plan de Calidad de Software para ChibchaWeb, tiene por objetivo
definir los procesos, métricas, responsabilidades y herramientas que aseguren la calidad
continua del sistema durante todas sus fases: desarrollo, implementación y operación. Está
diseñado para:

\begin{itemize}
  \item Garantizar la confiabilidad y disponibilidad del servicio, conforme al estándar del
        99,9\,\% de uptime ofrecido.
  \item Asegurar la integridad y protección de los datos mediante la correcta gestión de
        certificados SSL y respaldos automáticos.
  \item Mantener altos niveles de soporte y satisfacción al usuario, alineados con la
        experiencia técnica 24/7 proporcionada.
\end{itemize}

Este plan se aplicará a todas las áreas del sistema ChibchaWeb, incluyendo gestión de
dominios, administración de planes, facturación, interfaces de usuario multilingües (español,
inglés, portugués) y canales de soporte (centro de ayuda, tickets, soporte vía email o teléfono).

% ============================================================================
\section{Objetivos de calidad}
% ============================================================================

El presente plan de calidad tiene como propósito asegurar que ChibchaWeb cumpla con los
más altos estándares de calidad en todos los aspectos relevantes del sistema. A continuación,
se describen los objetivos específicos:

\subsection{Seguridad}

\begin{itemize}
  \item Garantizar la protección de los datos del cliente mediante el uso de certificados SSL,
        mecanismos de autenticación y control de acceso.
  \item Prevenir accesos no autorizados y asegurar la integridad de los respaldos automáticos.
\end{itemize}

\subsection{Rendimiento y calidad}

\begin{itemize}
  \item Asegurar tiempos de respuesta rápidos para las interfaces del usuario ($< 2$ segundos
        en operaciones comunes).
  \item Mantener una disponibilidad mínima del 99,9\,\%, tal como se promociona en la
        plataforma, incluso durante tareas de mantenimiento.
\end{itemize}

\subsection{Mantenibilidad}

\begin{itemize}
  \item Facilitar la incorporación de nuevas funcionalidades sin afectar la estabilidad
        existente.
  \item Asegurar que el código esté documentado, modularizado y probado adecuadamente
        para permitir su evolución futura.
\end{itemize}

\subsection{Usabilidad}

\begin{itemize}
  \item Ofrecer una experiencia de usuario fluida, clara y accesible, adaptable a distintos
        dispositivos.
  \item Mantener consistencia en la navegación entre módulos, como el panel de control,
        gestión de planes, dominios y soporte.
\end{itemize}

\subsection{Confiabilidad}

\begin{itemize}
  \item Asegurar que todas las funcionalidades del sistema se ejecuten correctamente, sin
        errores en condiciones normales.
  \item Disminuir la tasa de fallos críticos mediante un sistema robusto de pruebas
        automatizadas y revisiones de calidad.
\end{itemize}

% ============================================================================
\section{Estándares, métricas y normas aplicadas}
% ============================================================================

Esta sección define los estándares de desarrollo, normas de calidad y métricas que se han
usado y se emplearán para evaluar y controlar la calidad del sistema ChibchaWeb a lo largo
de su ciclo de vida.

\subsection{Estándares de desarrollo}

Para asegurar la coherencia, mantenibilidad y seguridad del código fuente de ChibchaWeb, se
aplican los siguientes estándares:

\begin{itemize}
  \item \textbf{PEP8} (Python Enhancement Proposal 8): para el formateo y estilo del código Python.
  \item \textbf{HTML5 y CSS3}: para garantizar accesibilidad, compatibilidad entre navegadores y
        estándares actuales del frontend.
  \item \textbf{Control de versiones con Git}: uso de ramas (branching), commits semánticos y
        revisión de código por pares.
\end{itemize}

\subsection{Normas y referencias de calidad}

\begin{itemize}
  \item \textbf{ISO/IEC 25010}: para la evaluación de la calidad del software con base en atributos
        como funcionalidad, confiabilidad, usabilidad, eficiencia, mantenibilidad y
        portabilidad.
  \item \textbf{IEEE 829 / ISO 29119}: para la planificación y documentación estructurada de
        pruebas.
  \item \textbf{ISO/IEC 27001}: lineamientos generales para la seguridad de la información,
        aplicados a nivel de hosting, respaldos y cifrado de datos.
\end{itemize}

\subsection{Métricas de calidad}

Estas métricas permitirán cuantificar y supervisar el nivel de calidad alcanzado durante y
después del desarrollo:

\begin{center}
\rowcolors{2}{white}{secbg}
\begin{tabularx}{\textwidth}{>{\raggedright}p{3.5cm} X >{\centering\arraybackslash}p{2.5cm}}
\toprule
\textbf{Métrica} & \textbf{Descripción} & \textbf{Umbral aceptable} \\
\midrule
Cobertura de pruebas unitarias
  & Porcentaje del código cubierto por pruebas automáticas
  & $\geq$ 80\,\% \\
Tasa de errores por módulo
  & Número de errores encontrados en cada módulo
  & $\leq$ 3 por módulo \\
Tiempo medio de resolución de errores
  & Tiempo promedio en corregir bugs desde su detección
  & $<$ 24 horas (críticos) \\
Tiempo de respuesta del servidor
  & Promedio de latencia por solicitud HTTP
  & $\leq$ 2 segundos \\
Porcentaje de disponibilidad
  & Tiempo activo de servicio / tiempo total
  & $\geq$ 99,9\,\% \\
Satisfacción del usuario (encuestas)
  & Valoración promedio en encuestas post-uso
  & $\geq$ 4.0 / 5.0 \\
\bottomrule
\end{tabularx}
\end{center}

% ============================================================================
\section{Herramientas y procedimientos de evaluación}
% ============================================================================

Para asegurar la medición y evaluación efectiva de la calidad del proyecto ChibchaWeb, se
utilizarán diversas herramientas y procedimientos:

\subsection{Herramientas de evaluación}

\begin{enumerate}
  \item \textbf{Google PageSpeed:} Para la medición de tiempo de respuesta.
  \item \textbf{Auditorías de Seguridad:} Se realizan a intervalos semestrales utilizando
        herramientas especializadas para garantizar la integridad y confidencialidad de los
        datos como HostedScan. Una vez lanzada la aplicación se comenzará con la primera
        medición.
  \item \textbf{Encuestas Anónimas:} Para medir y evaluar la satisfacción del usuario una vez
        realizado el lanzamiento del sistema.
\end{enumerate}

\subsection{Procedimientos de evaluación}

\begin{enumerate}
  \item \textbf{Monitoreo Continuo:} Las métricas de tiempo de respuesta y disponibilidad se
        medirán y registrarán de forma continua, con informes diarios y semanales
        respectivamente.
  \item \textbf{Auditorías de Seguridad:} Se llevarán a cabo auditorías exhaustivas cada seis meses
        para evaluar y mejorar la seguridad de los datos.
  \item \textbf{Encuestas de Satisfacción:} Trimestralmente, se realizarán encuestas anónimas a los
        usuarios para evaluar su satisfacción con el sistema.
\end{enumerate}

% ============================================================================
\section{Roles y responsabilidades}
% ============================================================================

Para garantizar la calidad del software durante el ciclo de vida del proyecto ChibchaWeb, se
han definido los siguientes roles, cada uno con responsabilidades específicas relacionadas con
la validación, aseguramiento y mejora continua de la calidad del sistema:

\subsection{Gerente del proyecto / Líder técnico}

\begin{itemize}
  \item Define y supervisa la ejecución del plan de calidad.
  \item Asegura que los recursos y tiempos destinados al aseguramiento de calidad sean
        respetados.
  \item Evalúa el cumplimiento de los estándares de desarrollo y los objetivos de calidad
        definidos.
\end{itemize}

\subsection{Desarrolladores}

\begin{itemize}
  \item Implementan las funcionalidades del sistema cumpliendo con las buenas prácticas de
        programación (PEP8, modularidad, control de versiones).
  \item Ejecutan pruebas unitarias y depuran errores identificados durante el desarrollo.
  \item Documentan el código y colaboran en la revisión de código entre pares.
\end{itemize}

\subsection{Desarrollador Tester / Analista de calidad}

\begin{itemize}
  \item Diseña y ejecuta casos de prueba para validar la funcionalidad, rendimiento y
        seguridad del sistema.
  \item Reporta errores encontrados y hace seguimiento hasta su resolución.
  \item Evalúa la experiencia de usuario y propone mejoras de usabilidad.
\end{itemize}

\subsection{Desarrollador responsable de soporte y atención al cliente}

\begin{itemize}
  \item Recibe retroalimentación de los usuarios finales en relación con problemas o
        sugerencias.
  \item Clasifica y canaliza los incidentes técnicos al equipo correspondiente.
  \item Participa en el análisis de satisfacción del cliente, contribuyendo a la mejora continua.
\end{itemize}

% ============================================================================
\section{Plan de acción correctiva}
% ============================================================================

En caso de identificar desviaciones o problemas en las métricas de calidad, se implementará
un plan de acción correctiva. Este plan incluirá:

\begin{enumerate}
  \item \textbf{Identificación y Análisis:}
        El equipo de desarrollo y los responsables de la seguridad y atención al cliente
        identificarán y analizarán las causas raíz de las desviaciones.

  \item \textbf{Implementación de soluciones:}
        Una vez identificadas las causas, se implementarán acciones correctivas específicas
        para abordar y resolver los problemas.

  \item \textbf{Monitoreo y Evaluación Continua:}
        Se llevará a cabo un seguimiento constante para asegurar que las acciones correctivas
        hayan solucionado los problemas detectados. Se realizarán ajustes según sea necesario
        para mantener los estándares de calidad establecidos.
\end{enumerate}

\end{document}
